% !PRINTED 362
\weberpage{362}
%<<CONT
but also the pair $x_3 y_3$, and consequently all the pairs $x_i y_i$, must occur, so that this complex is not different at all from the one derived from $x_1 y_1$. We therefore have the theorem:

\gesperrt{2. The pairs of a \gesperrt{Steiner} complex have the property that the eight points of contact of any two of these pairs lie on a conic, and that one always obtains the same complex, from whichever of the six pairs one may start.}

From this it follows that three double tangents of a \gesperrt{Steiner} complex, such as $x_1, y_1, x_2$, of which two form a pair of the complex, have their six points of contact on a conic. On the other hand, there are again systems of three double tangents (triples) whose six points of contact do not lie on a conic.

According to a designation introduced by \gesperrt{Frobenius}, three double tangents, such as $x_1, y_1, x_2$, whose points of contact lie on a conic, form a \gesperrt{syzygetic triple}. Three double tangents whose six points of contact do not lie on a conic form an \gesperrt{azygetic triple}. Correspondingly, we shall call four double tangents whose eight points of contact lie on a conic a \gesperrt{syzygetic quadruple}, and any system of double tangents of which every three are azygetic an \gesperrt{azygetic system}.

Here now the following important theorem holds:

\gesperrt{3. Three double tangents which occur in a \gesperrt{Steiner} complex in such a way that no two of them form a pair, such as $x_1, x_2, x_3$, are always azygetic.}

The proof of this theorem follows simply from the form of the equation (9). From it one sees first of all that the three lines $x_1, x_2, x_3$ do not intersect in one point; for such a point of intersection would lie on the curve $f$ and would therefore be a singular point. We can therefore introduce $x_1, x_2, x_3$ as coordinates and accordingly
\[
\begin{gathered}
y_1 = \alpha_1 x_1 + \alpha_2 x_2 + \alpha_3 x_3\\
y_2 = \beta_1 x_1 + \beta_2 x_2 + \beta_3 x_3\\
y_3 = \gamma_1 x_1 + \gamma_2 x_2 + \gamma_3 x_3
\end{gathered}
\]
%CONT>>
